\begin{Example} \label{verma_prim}
Consider $p=1$ case. Singular vector in $V^{\HVir}(h,0) $ is $u'_{1}=\big( L(-1) +\frac{h}{c_{L,I}}I(-1) \big) v_{h,0}$, and $u'_{1}$ generates a copy of $V^{\HVir}(h+1,0) $.

$r=1$: $\Psi \colon V^{W(2,2)}(0,0) \rightarrow V^{\HVir}(0,0) $ maps a singular vector $u'_1 = W(-1) v_{0,0}$ to $0$ and a~cosingular vector $u_1 = L(-1) v_{0,0}$ to $\HVir$-singular vector $v_1^- = L(-1) v_{0,0}$. We get the short exact sequence of $W(2,2)$-modules
\begin{gather*}
0\rightarrow L^{W(2,2)}(0,0) \rightarrow L^{\HVir}(0,0) \rightarrow L^{W(2,2)}(1,0) \rightarrow 0,
\end{gather*}
which is an expansion of (\ref{Psi}) considered as a homomorphism of $W(2,2)$-modules. The rightmost module is generated by a projective image of~$I(-1)v_{0,0}$. Therefore, $L^{\HVir}(c_L, c_{L,I})$ is generated over $W(2,2)$ by $v_{0,0}$ and $I(-1)v_{0,0}$.

$r\in \N$: In general, a cosingular vector $u_{rp}\in V^{W(2,2)}\left( \frac{1-r}{2},0\right) $ maps to a singular vector $v_r^- \in V^{
\HVir}\big( \frac{1-r}{2},0\big)$ of weight ${\frac{1+r}{2}}$.
\begin{gather*}
v_r^- = \prod\limits_{i=0}^{r-1}\left( L(-1) +\frac{1-r+2i}{2c_{L,I}}I(-1) \right) v_{\frac{1-r}{2},0}.
\end{gather*}
\end{Example}

\section[Screening operators and $W(2,2)$-algebra]{Screening operators and $\boldsymbol{W(2,2)}$-algebra}\label{screenings}

We think that the vertex algebra $L^{W(2,2)} (c_L, c_W) $ is a very interesting example of non-rational vertex algebra, which admits similar fusion ring of representations as some ${\mathcal W}$-algebras appearing in LCFT (cf.\ \cite{A1,AdM1,CM,FGST}). Since $\mathcal W$-algebras appearing in LCFT are realized as kernels of screening operators acting on certain modules for Heisenberg vertex algebras, it is natural to ask if $L^{W(2,2)} (c_L, c_W) $ admits similar realization. In~\cite{nas} we embedded the $W(2,2)$-algebra as a subalgebra of the Heisenberg--Virasoro vertex algebra. In this section we shall construct a~screening operator $S_1$ such that the kernel of this operator is exactly $L^{W(2,2)} (c_L, c_W) $.

Let us f\/irst construct a non-semisimple extension of the vertex algebra $L^{\HVir} (c_L, c_{L,I})$. Recall that the Lie algebra $\HVir$ admits the triangular decomposition~(\ref{triangular}). Let $E=\operatorname{span}_{\C} \{ v^{0}, v^{1} \}$ be $2$-dimensional $\HVir^{\ge 0} = \HVir^{0} \oplus \HVir^{+}$-module such that $\HVir ^+ $ acts trivially on $E$ and
\begin{gather*}
L(0) v^{i} = v^{i} , \qquad i=0,1, \qquad I(0) v^1 = v^0, \qquad I(0) v^0 = 0, \\
C_L v^{i} = c_L v^{i}, \qquad C_{L,I} v^{i} = c_{L,I} v^{i}, \qquad C_I v^{i} = 0, \qquad i=1,2.
\end{gather*}
Consider now induced $\HVir$-module
\begin{gather*}
\widetilde E = U(\HVir) \otimes _{U (\HVir ^{\ge 0} )} E.
\end{gather*}
By construction, $\widetilde E$ is a non-split self-extension of the Verma module $V^{\HVir} (1,0)$:
\begin{gather*}
0 \rightarrow V^{\HVir} (1,0) \rightarrow \widetilde E \rightarrow V^{\HVir} (1,0 ) \rightarrow 0.
\end{gather*}
Moreover, $\widetilde E$ is a restricted module for $\HVir$ and therefore it is a module over vertex operator algebra $L^{\HVir} (c_L, c_{L,I})$. Since
\begin{gather*}
\widetilde E \cong E \otimes U(\HVir^{-})
\end{gather*}
as a vector space, the operator $L(0)$ def\/ines a $\Z_{\ge 0}$-gradation on $\widetilde E$.

Note that $ (L(-1) + I(-1) / c_{L,I} ) v_0$ is a singular vector in $\widetilde E$ and it generates the proper submodule. Finally we def\/ine the quotient module
\begin{gather*}
\mathcal U = \frac{ \widetilde E} {U(\HVir) . (L(-1) + I(-1) / c_{L,I} ) v_0}.
\end{gather*}

\begin{Proposition}
$\mathcal U$ is a $\Zp$-graded module for the vertex operator algebra $L^{\HVir} (c_L, c_{L,I} )$:
\begin{gather*}
\mathcal U = \bigoplus _{m \in {\Zp} } \mathcal U (m), \qquad L(0) \vert \mathcal U(m) \equiv (m+1) \operatorname{Id}.
\end{gather*}
The lowest component $\mathcal U (0) \cong E$. Moreover, $\mathcal U$ is a non-split extension of the Verma module $V^{\HVir} (1,0)$ by the simple highest weight module $L^{\HVir} (1,0)$:
 \begin{gather*}
 0 \rightarrow L^{\HVir} (1,0) \rightarrow \mathcal U \rightarrow V^{\HVir} (1,0 ) \rightarrow 0.
 \end{gather*}
 \end{Proposition}
 \begin{proof}
By construction $\mathcal U $ is a graded quotient of a $\Z_{\ge 0}$-graded $L^{\HVir}(c_L, c_{L,I})$-module $\widetilde E$. The lowest component is $\mathcal U (0) \cong E$. Submodule $U(\HVir).v^0$ is isomorphic to $L^{\HVir}(1,0)$, and the projective image of $v^1$ generates the Verma module $V^{\HVir}(1,0)$ since $I(0)v^1=v^0$. For the same reason, this exact sequence does not split.
\end{proof}

Now we consider $L^{\HVir}(c_L, c_{L,I} )$-module
\begin{gather*}
\mathcal V^{\rm ext}:= L^{\HVir}(c_L, c_{L,I} ) \oplus \mathcal U.
\end{gather*}
By using \cite[Theorem 4.8.1]{LL} (see also \cite{AdM-selecta,Li1}) we have that $ \mathcal V^{\rm ext}$ has the structure of a~vertex operator algebra with vertex operator map $Y_{\rm ext}$ def\/ined as follows:{\samepage
\begin{gather*}
Y_{\rm ext} ( a_1 + w_1 , z ) (a_2 + w_2) = Y(a_1, z) (a_2+w_2) + e^{z L(-1)} Y (a_2 ,-z ) w_1,
\end{gather*}
where $a_1, a_2 \in L^{\HVir}(c_L, c_{L,I} )$, $ w_1, w_2 \in \mathcal U$.}

Take now $v^ i \in E \subset \mathcal U$, $i=0,1$ as above and def\/ine
\begin{gather*}
S_i (z) = Y_{\rm ext} \big(v^i, z\big) = \sum _{n \in \Z} S_i (n) z^{-n-1}.
\end{gather*}
By construction
\begin{gather*}
S_1 (z) \in \operatorname{End} \big( L^{\HVir} (c_L,c_{L,I}), L^{\HVir}(1,0) \big) ((z)).
\end{gather*}

\begin{Proposition}
For all $n,m\in \Z$ we have:
\begin{gather*}
[L(n), S_i (m)] = - m S_i (n+m), \qquad i=0,1,\\
[W(n), S_0 (m) ] = 0, \qquad [W(n), S_1(m) ] = 2 m c_{L,I} S_0( n+m).
\end{gather*}
In particular, $S_0 (0) $ and $S_1(0)$ are screening operators. Moreover,
\begin{gather*}
S_1= S_1(0) \colon \ L^{\HVir} (c_L, c_{L,I}) \rightarrow L^{\HVir} (1,0)
\end{gather*}
is nontrivial and $ S_1 (0) I(-1) {\bf 1} = - v_0. $
\end{Proposition}

\begin{proof} Since $L(k) v^{i} = \delta_{k,0} v^{i}$ for $k \ge 0$, commutator formula gives that
\begin{gather*}
[L(n), S_i(m)] = - m S_i (n+m).
\end{gather*}
Next we calculate $[W(n), S_1 (m)]$. We have
\begin{gather*}
W(-1) v^1 = 2 I(-1) v^0 = - 2 c_{L,I} L(-1) v^0, \\
W(0) v^1 = - 2 c_{L,I} v^ 0, \qquad W(n) v^1 = 0, \qquad n \ge 0.
\end{gather*}
This implies that
\begin{gather*}
[W(n), S_1 (m) ] = 2 c_{L,I} m S_0 (n+m).
\end{gather*}
Since $W(n) v^0 = 0$ for $n \ge -1$ we get
\begin{gather*}
[ W(n), S_0 (m) ] = 0.
\end{gather*}
Therefore we have proved that $S_i(0)$, $i= 0,1$ are screening operators. Next we have
\begin{gather*}
S_1(0) I(-1) {\bf 1} = \operatorname{Res}_z Y_{\rm ext}\big(v^1, z\big) I(-1) = \operatorname{Res}_z e ^{z L(-1)} Y( I(-1) {\bf 1}
,- z) v^1 = - v_0.
\end{gather*}
The proof follows.
\end{proof}

\begin{Theorem}
$S_1$ is a derivation of the vertex algebra $\mathcal V^{\rm ext}$ and we have
\begin{gather*}
\Ker_{ L^{\HVir} (c_L,c_{L,I}) } S_1 \cong L ^{W(2,2)} (c_L, c_W ).
\end{gather*}
\end{Theorem}
\begin{proof}
By construction $S_1 = \operatorname{Res}_z Y_{\rm ext} (v^1, z)$, so $S_1$ is a derivation so $\overline W = \Ker_{ L^{\HVir} (c_L,c_{L,I}) } S_1$ is a vertex subalgebra of $L^{\HVir} (c_L,c_{L,I})$. Since
\begin{gather*}
S_1 L(-2) {\bf 1} = S_1 W(-2) {\bf 1} = 0
\end{gather*}
we have that $L ^{W(2,2)} (c_L, c_W ) \subset \overline W$. Since $S_1 I (-1) {\bf 1} \ne 0$, we have that $I(-1) {\bf 1}$ does not belong to $\overline W$. By using the fact that $L^{\HVir} (c_L,c_{L,I})$ is as $W(2,2)$-module generated by singular vector $\bf{1}$ and cosingular vector $I(-1) {\bf 1}$ (see Example~\ref{verma_prim}) we get that $\overline W = L ^{W(2,2)} (c_L, c_W )$. The proof follows.
\end{proof}
